% Contenu adapté pour : Samuel SIKATI KENMOGNE -- version Ingénieur QA / Test logiciel
%
% Template license:
% CC BY-SA 4.0 (https://creativecommons.org/licenses/by-sa/4.0/)

% Commandes de mise en forme définies dans yaac-another-awesome-cv.cls :
%   \company{<url>}{<nom>}   entreprise avec lien (petites capitales, pointillé)
%   \companynolink{<nom>}    entreprise sans site officiel
%   \projet{<intitulé>}      intitulé de projet / contexte (poids medium)

%Section: Expérience QA
\sectionTitle{Expériences Professionnelles}{\faSuitcase}
\begin{experiences}
  \experience[CDD]
    {Aujourd'hui}   {Ingénieur QA / Test Logiciel}{\company{https://www.orange.cm}{Orange Cameroun}}{Douala}
    {Nov. 2025} {
                      \projet{Qualification de la migration d'une plateforme mobile money}
                      \begin{itemize}
                        \item Contribution à la stratégie de test et planification de chantiers de recette
                        \item Conception de \textbf{1\,547 cas de test} répartis en cahiers de recette couvrant tous les chantiers de la migration
                        \item Tests d'intégration système et end-to-end multicanaux (USSD, mobile, web) des flux transactionnels
                        \item Qualification de la bascule : migration de données, plan de reprise d'activité (PRA), performance (Locust)
                        \item Automatisation de \textbf{640 tests de régression} API (Bruno) et web (Playwright), exécutés en CI (Azure DevOps)
                        \item Animation des campagnes UAT avec les métiers et qualification de \textbf{262 anomalies} avant la bascule
                      \end{itemize}
                      \projet{Recette fonctionnelle d'une solution de gestion du courrier entrant (Maarch Courrier)}
                      \begin{itemize}
                        \item Paramétrage, validation des exigences et tests d'acceptation ; formation et support au déploiement
                      \end{itemize}
                    }
                    {Stratégie de test,UAT,Tests E2E,Tests API,Régression,PRA,Playwright,Bruno,Locust,Azure DevOps}
  \emptySeparator
  \experience[Stage]
    {Juin 2025}   {Ingénieur Logiciel Junior}{\company{https://societegenerale.cm/fr/}{Société Générale Cameroun}}{Douala}
    {Jan. 2025} {
                      \projet{Développement d'une solution de suivi des anomalies opérationnelles bancaires}
                      \begin{itemize}
                        \item Modélisation et implémentation du workflow de plans d'action correctifs (échéances, relances, preuves)
                        \item Mise en place de tableaux de bord KPI et automatisation du reporting, jusque-là manuel
                        \item Réalisation des tests end-to-end de l'application avant sa mise en production
                        \item \textbf{Réduction de 25\,\% du temps de traitement} des anomalies et \textbf{de 35\,\% des erreurs manuelles}
                      \end{itemize}
                    }
                    {.NET Core,C\#,Tests E2E,Gestion des anomalies,Traçabilité,KPI,Tableaux de bord,Banque}
  \emptySeparator
  \experience[Stage]
    {Fév. 2024}   {Développeur Logiciel}{\company{https://www.boissonsducameroun.com}{Boissons du Cameroun}}{Douala}
    {Sept. 2023} {
                      \projet{Digitalisation du processus d'analyse des causes profondes des incidents (SCAT -- DNV)}
                      \begin{itemize}
                        \item Modélisation de la matrice SCAT (contacts, causes immédiates et fondamentales, actions de contrôle)
                        \item Développement d'une application web (Symfony, Flask) dédiée aux investigations d'incidents
                        \item \textbf{Réduction de 40\,\% du temps d'investigation} et amélioration de la qualité des rapports d'analyse
                      \end{itemize}
                    }
                    {Symfony,PHP,Flask,Python,Application web,SCAT -- DNV,HSE,Analyse des causes profondes}
  \emptySeparator
  \experience[Stage]
    {Avr. 2023}   {Analyste Programmeur}{\company{https://www.perenco.com}{Perenco Cameroun}}{Douala}
    {Jan. 2023} {
                      \projet{Optimisation du contrôle qualité QHSE par une application mobile et des macros Excel}
                      \begin{itemize}
                        \item Développement d'une application mobile Android (Kotlin) et automatisation des contrôles en VBA Excel
                        \item \textbf{Amélioration de 20\,\% de l'efficacité des audits} et meilleure conformité réglementaire
                      \end{itemize}
                    }
                    {Kotlin,Android,Application mobile,VBA,Excel,QHSE,Audit,Conformité réglementaire,Oil \& Gas}
  \emptySeparator
  \experience[Stage]
    {Oct. 2022}   {Développeur Back-End}{\companynolink{WinMoney Cameroun}}{Yaoundé}
    {Juin 2022} {
                      \projet{Développement de l'API RESTful d'une plateforme de création de contenus visuels}
                      \begin{itemize}
                        \item Conception et implémentation des endpoints REST et de la logique métier avec Spring Boot
                        \item Support de \textbf{jusqu'à 500 requêtes/minute} avec un temps de réponse moyen de 300\,ms
                      \end{itemize}
                    }
                    {Java,Spring Boot,API REST,Back-end,Performance,Scalabilité}
  \emptySeparator
  \experience[Stage]
    {Oct. 2021}   {Développeur Back-End Junior}{\company{https://kiama.cm}{Kiama S.A}}{Yaoundé}
    {Juin 2021} {
                      \projet{Mise en place d'un système de recommandation multicanal (email/SMS) pour \company{https://ekiosque.cm}{ekiosque.cm}}
                      \begin{itemize}
                        \item Personnalisation des recommandations par machine learning (PHP-ML) et diffusion par email et SMS
                        \item \textbf{Augmentation de 15\,\% de l'engagement utilisateur} et amélioration du taux de conversion
                      \end{itemize}
                    }
                    {Symfony,PHP,PHP-ML,Machine Learning,Recommandation,Email,SMS,Presse numérique}
\end{experiences}
